\Needspace{14\baselineskip}
\section{Final statements and their trust boundary}\label{app:statements}
The final types below are the output of Lean's \code{\#print}, with
line wrapping only and proof terms omitted. The type of \code{Point}
is $\mathbb R\times\mathbb R$.
\begin{lstlisting}
theorem ES7FinalVerified.es7_le_81 :
  (*@$\forall$@*) (S : Finset ES8Geometry.Point),
  ES7Final.GeneralPosition S (*@$\to$@*) 81 (*@$\le$@*) S.card (*@$\to$@*)
  ES7Geometry.HasConvexHeptagon (*@$\uparrow$@*)S

theorem ES8Final.es8_le_319 :
  (*@$\forall$@*) (S : Finset ES8Geometry.Point),
  ES8Final.GeneralPosition S (*@$\to$@*) 319 (*@$\le$@*) S.card (*@$\to$@*)
  ES8Geometry.HasConvexOctagon (*@$\uparrow$@*)S
\end{lstlisting}
Both general-position definitions unfold to the following expression:
\begin{lstlisting}
fun S => (*@$\forall$@*) a (*@$\in$@*) S, (*@$\forall$@*) b (*@$\in$@*) S, (*@$\forall$@*) c (*@$\in$@*) S,
  a (*@$\ne$@*) b (*@$\to$@*) b (*@$\ne$@*) c (*@$\to$@*) a (*@$\ne$@*) c (*@$\to$@*) ES8Geometry.orient a b c (*@$\ne$@*) 0
\end{lstlisting}
Here \code{orient p q r} is
$(q_x-p_x)(r_y-p_y)-(q_y-p_y)(r_x-p_x)$. The two polygon predicates
unfold respectively to
\begin{lstlisting}
fun S => (*@$\exists$@*) T, T.card = 7 (*@$\land$@*) (*@$\uparrow$@*)T (*@$\subseteq$@*) S (*@$\land$@*) ES8Geometry.ConvexPosition (*@$\uparrow$@*)T
fun S => (*@$\exists$@*) T, T.card = 8 (*@$\land$@*) (*@$\uparrow$@*)T (*@$\subseteq$@*) S (*@$\land$@*) ES8Geometry.ConvexPosition (*@$\uparrow$@*)T
\end{lstlisting}
The underlying definition and the new bridge are
\begin{lstlisting}
def ConvexPosition (S : Set Point) : Prop :=
  (*@$\forall$@*) p (*@$\in$@*) S, p (*@$\notin$@*) convexHull (*@$\mathbb R$@*) (S \ {p})

theorem convexPosition_iff_convexIndependent (S : Set Point) :
  ConvexPosition S (*@$\leftrightarrow$@*) ConvexIndependent (*@$\mathbb R$@*) (((*@$\uparrow$@*)) : S (*@$\to$@*) Point) := by
  exact convexIndependent_set_iff_notMem_convexHull_sdiff.symm
\end{lstlisting}
The bridge uses only the three standard axioms listed below. The
companion module applies it to both endpoints. Thus their convexity
conclusions can be stated using mathlib's standard predicate without
strengthening a hypothesis or trusting another finite evaluation.

\Needspace{10\baselineskip}
\begingroup
\lstset{aboveskip=3pt,belowskip=3pt}
\subsection{The restricted theorem and its chain predicates}
The restricted upper bound uses the following declaration and definitions.
The upper bound is formalized; the lower bound is supplied by the cited
construction and the independently checked 25-point example.
\begin{lstlisting}
theorem ES7CapacityStage.restricted25 :
  ES7Residual.RestrictedBound 25
\end{lstlisting}
\Needspace{8\baselineskip}
\begin{lstlisting}
def ES7Residual.RestrictedBound (K : (*@$\mathbb N$@*)) : Prop :=
  (*@$\forall$@*) {n} (c : ES8BalancedGeometry.Configuration n),
    ((*@$\neg$@*)(colorConfig c.cup true).HasNCap 5 (Finset.range n)) (*@$\to$@*)
    ((*@$\neg$@*)(colorConfig c.cup true).HasNCup 6 (Finset.range n)) (*@$\to$@*)
    ((*@$\neg$@*)HasConvexHeptagon (Set.range c.point)) (*@$\to$@*) n (*@$\le$@*) K
\end{lstlisting}
\Needspace{8\baselineskip}
\begin{lstlisting}
structure ES8BalancedGeometry.Configuration (n : (*@$\mathbb N$@*)) where
  point : Fin n (*@$\to$@*) Point
  xOrder : StrictMono (fun i => (point i).1)
  generalPosition : (*@$\forall$@*) u v w : Fin n,
    u < v (*@$\to$@*) v < w (*@$\to$@*) orient (point u) (point v) (point w) (*@$\ne$@*) 0
\end{lstlisting}
For valid indices, \code{c.cup u v w} is true exactly when the
orientation determinant is positive; out-of-range indices return false.
The notation \code{colorConfig} refers to \code{ES8ActualRanks.colorConfig}:
\begin{lstlisting}
def colorConfig (cup : (*@$\mathbb N$@*) (*@$\to$@*) (*@$\mathbb N$@*) (*@$\to$@*) (*@$\mathbb N$@*) (*@$\to$@*) Bool) (color : Bool) : Config (*@$\mathbb N$@*) :=
  { Cup3 := fun a b c => cup a b c = color,
    DecidableCup3 := inferInstance }
\end{lstlisting}
In the Baek port's \code{Config} namespace, with a linearly ordered
vertex type $\alpha$ and \code{C : Config} $\alpha$, the predicates are
\begin{lstlisting}
def Cap3 (a b c : (*@$\alpha$@*)) : Prop := (*@$\neg$@*)C.Cup3 a b c
def Cap (l : List (*@$\alpha$@*)) : Prop :=
  l.IsChain ((*@$\cdot$@*) < (*@$\cdot$@*)) (*@$\land$@*) l.Chain3' C.Cap3
def Cup (l : List (*@$\alpha$@*)) : Prop :=
  l.IsChain ((*@$\cdot$@*) < (*@$\cdot$@*)) (*@$\land$@*) l.Chain3' C.Cup3
def NCap (n : (*@$\mathbb N$@*)) (l : List (*@$\alpha$@*)) : Prop :=
  C.Cap l (*@$\land$@*) l.length = n
def NCup (n : (*@$\mathbb N$@*)) (l : List (*@$\alpha$@*)) : Prop :=
  C.Cup l (*@$\land$@*) l.length = n
def HasNCap (n : (*@$\mathbb N$@*)) (S : Finset (*@$\alpha$@*)) : Prop :=
  (*@$\exists$@*) l : List (*@$\alpha$@*), C.NCap n l (*@$\land$@*) l.In S
def HasNCup (n : (*@$\mathbb N$@*)) (S : Finset (*@$\alpha$@*)) : Prop :=
  (*@$\exists$@*) l : List (*@$\alpha$@*), C.NCup n l (*@$\land$@*) l.In S
\end{lstlisting}
Here \code{IsChain} tests consecutive order comparisons,
\code{Chain3'} tests consecutive triples, and \code{l.In S} means
that every listed vertex lies in $S$. Strict order ensures distinct
vertices. For actual points these are precisely the slope definitions
of cups and caps in Section~\ref{sec:geometry}. The archived statement
audit prints the fully qualified definitions directly from Lean.
\endgroup

\clearpage
\subsection{Exact final axiom reports}
The following is the complete \code{\#print axioms} output for the
two original endpoints. The 19 and eight names ending in
\code{native\_decide.ax\_1\_1} are the disclosed native checks.
\lstinputlisting{axiom_output.txt}
